\begin{center}{\zihao{4}\heiti 摘\quad 要}\end{center}
\addcontentsline{toc}{section}{摘要}

针对普通视频中微小振动难以辨识、运动位置与频率信息难以对应的问题，本作品设计了基于AMD Zynq-7000的相位运动放大与区域微振检测系统。系统面向振动与信号处理教学、机械结构低频振动观察，在同一板端终端上集成运动可视化、区域位移分析和交互控制。

系统采用OV9281采集$1280\times800$图像，由可编程逻辑完成多尺度分解、局部相位处理、时间滤波和图像重建；独立测量支路从原始采集帧的固定$256\times256$区域估计位移，经256点滑动窗口和直接离散傅里叶变换生成129个单边频点。FPGA以流水线、行缓存和专用运算模块承载密集计算，处理系统负责相机配置、帧调度、显示提交及串口屏交互，实现采集、计算、显示与控制的协同运行。

扬声器双频实板实验在5~Hz与10~Hz激励附近得到对应谱峰。v19补充记录完成三次断电冷启动、按键及触屏交互确认；两组黄金参考与HLS C模型比较中，已归档样例的最大绝对误差不超过1个RAW10码值。v18版本在固定配置下连续运行32分钟，平均处理速率为60.37969~frame/s，显示速率为30.00000~frame/s。Vivado~2025.2基线实现使用LUT~138026个、FF~206397个、BRAM~744.5个Tile和DSP~826个，BRAM利用率为98.61\%，WNS为$+0.177$~ns。资源分摊结果将主要片上存储开销定位到运动放大和ROI模块，为缓存复用与后续扩展提供了依据。

本作品通过整帧运动增强与原始ROI测量双支路设计，将运动形态和局部主频在同一终端对应呈现；通过缓存所有权管理、参数生效确认及启停排空机制，协调连续视频处理与用户操作，形成可独立运行的视觉微振观察与分析终端。

\vspace{.8em}
\noindent{\zihao{4}\heiti 关键字}\quad 相位运动放大；Zynq；区域位移估计；振动频谱；软硬件协同
